%::::::::::::: LIBRERIAS    ::::::::::::::
%:::::::::::::::::::::::::::::::::::::::::

\usepackage[spanish, es-ucroman]{babel}
\usepackage[utf8]{inputenc}
\usepackage{fancyhdr}
\usepackage{graphicx}
\usepackage[inner=2.54cm,outer=2.54cm,top=2.54cm,bottom=2.54cm]{geometry}
\usepackage{newtxtext,newtxmath}
\usepackage{Estilo}
\let\openbox\relax
\usepackage{amsmath}
\usepackage{ragged2e}
\usepackage{amsthm}
\usepackage{etoolbox}
\usepackage{chngcntr}
\usepackage{setspace}  
\usepackage{multirow}
\let\Bbbk\relax
\usepackage{amssymb}
\usepackage{tikz}
\usepackage[breaklinks]{hyperref}
\usepackage{caption}
\usepackage{amsfonts}
\usepackage{natbib}
\usepackage{float}
\usepackage{titlesec}
\usepackage{titletoc}
\usepackage[T1]{fontenc}
\usepackage{chngcntr}
\usepackage{enumitem}
\usepackage{array}
\usepackage{pdflscape}
\usepackage{subcaption}
\usepackage{longtable}
\usepackage{tocloft}
\setlength{\bibhang}{1.27cm}
\raggedbottom

%:::::::::::::::::::::::::::::::::::::::::::::
%:::::::::::::  DECLARACIONES   ::::::::::::::
%:::::::::::::::::::::::::::::::::::::::::::::
\renewcommand{\thetable}{\arabic{table}}
\renewcommand{\thefigure}{\arabic{figure}}
\renewcommand{\cfttabpresnum}{\textbf{Tabla~}} 
\renewcommand{\cftfigpresnum}{\textbf{Figura~}}
\renewcommand{\cfttabaftersnum}{\hspace{1em}}
\renewcommand{\cfttabfont}{\normalfont} 
\renewcommand{\cfttabpagefont}{\normalfont}

% Para el título del trabajo
\newcommand{\mytitle}[1]{
    \begin{hyphenrules}{nohyphenation}
        #1
    \end{hyphenrules}
}

% Encabezado y Pie de Página
\pagestyle{fancy}
\fancyhf{}
\renewcommand{\headrulewidth}{0pt}
\fancyfoot[C]{\thepage}

% Formato de los capítulos
\titleformat{\chapter}[display]
  {\normalfont\bfseries\centering\large}
  {\MakeUppercase{\chaptertitlename}\ \Roman{chapter}}
  {0pt}
  {\MakeUppercase}
\titlespacing*{\chapter}{0pt}{0pt}{10pt}

% Formato de capítulo no numerado
\titleformat{name=\chapter,numberless}[display]
  {\normalfont\bfseries\large\centering}
  {}
  {0pt}
  {\MakeUppercase}

% Formato de secciones
\titleformat{\section}
  {\normalfont\large\bfseries}
  {\thesection}
  {1em}
  {}

% Formato de subsecciones
\titleformat{\subsection}
  {\normalfont\large\bfseries\itshape}
  {\bfseries\itshape\thesubsection}
  {1em}
  {}


% Nombre de los índices
\addto\captionsspanish{\renewcommand{\contentsname}{\begin{center} \large ÍNDICE \end{center}}}
\addto\captionsspanish{\renewcommand{\listtablename}{\begin{center} \large ÍNDICE DE TABLAS \end{center}}}
\addto\captionsspanish{\renewcommand{\listfigurename}{\begin{center} \large ÍNDICE DE FIGURAS \end{center}}}
\addto\captionsspanish{\renewcommand{\bibname}{\large{BIBLIOGRAFÍA}}}


% Formato de las tablas
\captionsetup[table]{name={Tabla},labelfont={bf,sl},textfont={sl},labelsep=newline,justification=raggedright,
singlelinecheck=false}
\counterwithout{table}{chapter}
\setlength{\cfttabnumwidth}{5.5em}
\setlength{\cfttabindent}{0em}

% Formato de las figuras
\captionsetup[figure]{labelfont={bf,sl},textfont={sl},labelsep=newline,justification=raggedright,
singlelinecheck=false}
\counterwithout{figure}{chapter}
\setlength{\cftfignumwidth}{4.5em}  
\setlength{\cftfigindent}{0em} 

%:::::::::::::::::::::::::::::::::::::::::
%:::::::::::  DEFINIR DEF, TEO Y COR :::::
%:::::::::::::::::::::::::::::::::::::::::
\newtheorem{definition}{Definición}[section]
\newtheorem{theorem}{Teorema}[section]
\newtheorem{corolary}{Corolario}[section]
% Placeholder para figuras pendientes
\newcommand{\figpend}[1]{%
  \begin{center}
    \fbox{\parbox{0.8\textwidth}{\centering\textit{[Figura pendiente: #1]}}}\par
  \end{center}
}
